\startSection
\selectlanguage{spanish}
\sectionTitle{Himno Nacional de Bolivia}

\begin{question}{En la historia musical boliviana, ¿quién compuso la música del Himno Nacional de Bolivia?}
  \choice{José Ignacio de Sanjinés}
  \choice{Eduardo Abaroa}
  \choice[correct]{Leopoldo Benedetto Vincenti}
  \choice{Franz Tamayo}
  \choice{Modesto Omiste}
\end{question}

\begin{question}{La letra del Himno Nacional de Bolivia expresa el espíritu patriótico de la etapa republicana. ¿Quién fue su autor?}
  \choice{Ricardo José Bustamante}
  \choice{Franz Tamayo}
  \choice{Adela Zamudio}
  \choice{Manuel José Cortés}
  \choice[correct]{José Ignacio de Sanjinés}
\end{question}

\begin{question}{¿En qué año fue estrenado oficialmente el Himno Nacional de Bolivia, en el marco de una celebración cívica?}
  \choice{1825}
  \choice{1842}
  \choice[correct]{1845}
  \choice{1851}
  \choice{1879}
\end{question}

\begin{question}{¿Cuál es el primer verso del Himno Nacional de Bolivia?}
  \choice{\enquote{De la Patria, el alto nombre}}
  \choice{\enquote{Compatriotas, la Patria muriendo}}
  \choice[correct]{\enquote{Bolivianos, el hado propicio}}
  \choice{\enquote{Gloria a Bolivia, tierra bendita}}
  \choice{\enquote{Surge ya la nueva aurora}}
\end{question}

\begin{question}{Además del coro, ¿cuántas estrofas contiene oficialmente el Himno Nacional de Bolivia?}
  \choice{Tres estrofas}
  \choice{Cuatro estrofas}
  \choice{Cinco estrofas}
  \choice{Seis estrofas}
  \choice[correct]{Nueve estrofas}
\end{question}

\begin{question}{En el himno, la oposición entre servidumbre y libertad cumple principalmente qué función retórica?}
  \choice{Describir un paisaje geográfico.}
  \choice{Narrar una anécdota personal del compositor.}
  \choice[correct]{Contrastar el pasado colonial con la afirmación republicana.}
  \choice{Enumerar las regiones administrativas del país.}
  \choice{Explicar una regla musical de interpretación coral.}
\end{question}
\shuffleSection
\renderSection